% lang/pl.tex
% -----------------------------------------------------------------------
% Las cadenas de texto de los cuadernos en polaco, las mismas que
% lang/es.tex, una a una (\booklang = pl, ver preamble.tex): «Pierwsze
% słowa» (slowa.tex) y, cuando lleguen, los otros niveles. Lo que cambia
% cada día (las palabras, los enunciados, la clave) lo escribe el
% generador de cada cuaderno.
%
% Lo que se le dice a la niña o al niño no tiene género ("umiesz",
% "samodzielnie", "dla tych, którzy..."): el cuaderno vale para los dos.
% -----------------------------------------------------------------------

% Total de días del cuaderno: el generador comprueba que coincide con el
% suyo (ver lang/es.tex).
\newcommand{\totaldias}{260}

\newcommand{\lblFecha}{Data}
\newcommand{\lblNotas}{Uwagi}
\newcommand{\lblDia}{Dzień}
\newcommand{\lblSemana}{Tydzień}
% "Część", como "Part" en inglés: cada parte del año es una estación.
\newcommand{\lblTrimestre}{Część}
\newcommand{\lblRespuesta}{Odpowiedź}

% "Czytanie", y no "Przeczytane": la línea de la cabecera (\cabeceraDia,
% preamble.tex) es la misma que en español, y con "samodzielnie" no cabía.
\newcommand{\lblLeido}{Czytanie}
\newcommand{\lblLeidoSola}{samodzielnie}
\newcommand{\lblLeidoConAyuda}{z pomocą}
\newcommand{\lblLeidoConDificultad}{z trudem}

% Una instrucción de lectura por trimestre (ver lang/es.tex).
\newcommand{\lblInstruccionLecturaUno}{Czytaj na głos, powoli, pokazując palcem każde słowo}
\newcommand{\lblInstruccionLecturaDos}{Przeczytaj na głos -- dwa razy}
\newcommand{\lblInstruccionLecturaTres}{Czytaj na głos, jak w teatrze}
\newcommand{\lblInstruccionLecturaCuatro}{Najpierw przeczytaj po cichu, dla siebie. Potem na głos}

\newcommand{\lblDibuja}{Narysuj}
\newcommand{\lblCompleta}{Dokończ}
\newcommand{\lblCopia}{Przepisz}
\newcommand{\lblResponde}{Odpowiedz}
\newcommand{\lblRelaciona}{Połącz}
\newcommand{\lblAdivina}{Zgadnij}
\newcommand{\lblCrea}{Wymyśl}
\newcommand{\lblRepasa}{Powtórka}
\newcommand{\lblRodea}{Zakreśl}
\newcommand{\lblVerdaderoFalso}{Prawda czy fałsz}
\newcommand{\lblBusca}{Znajdź}
\newcommand{\lblOrdena}{Ułóż po kolei}
\newcommand{\lblRelee}{Przeczytaj jeszcze raz}

\newcommand{\lblVerdadero}{P}
\newcommand{\lblFalso}{F}

\newcommand{\lblInstruccionCopia}{Przepisz dzisiejsze zdanie na linii.}
\newcommand{\lblInstruccionCompleta}{Zamień te linie w rysunek. Użyj wyobraźni!}
\newcommand{\lblInstruccionRodea}{Zakreśl dobrą odpowiedź.}
\newcommand{\lblInstruccionVerdaderoFalso}{Prawda czy fałsz? Zakreśl właściwą kratkę przy każdym zdaniu.}
\newcommand{\lblInstruccionOrdena}{Te zdarzenia są pomieszane. Wpisz w kratki 1, 2 i 3, w takiej kolejności, w jakiej się wydarzyły.}
\newcommand{\lblInstruccionRelee}{W tym tygodniu wszystko razem. Przeczytaj jeszcze raz, od początku do końca.}

% "Traza": la letra, mayúscula y minúscula, con su palabra -- "M jak
% mama", como en las cartillas polacas.
\newcommand{\lblTraza}{Pisz po śladzie}
\newcommand{\lblTrazoDe}{jak}
\newcommand{\lblInstruccionTrazo}{Poprowadź palcem albo ołówkiem po kropkach. Potem napisz literę samodzielnie na linii.}

\newcommand{\lblTituloPortada}{Uczę się czytać}
\newcommand{\lblSubtituloPortada}{Codziennie nowy rozdział}
\newcommand{\lblSubtituloPortadaMetodo}{Zeszyt do codziennego czytania, żeby czytać coraz płynniej -- a nie żeby uczyć się od zera}
\newcommand{\lblPortadaNombre}{Imię}
\newcommand{\lblPortadaCurso}{Klasa}

% -----------------------------------------------------------------------
% «Pierwsze słowa», el cuaderno de primeras palabras en polaco
% (slowa.tex, ver preamble-palabras.tex y tools/gen_palabras.py --libro
% slowa).
% -----------------------------------------------------------------------

% Instrucción de lectura (título de la caja verde), una por trimestre:
% sílaba a sílaba al principio, de corrido al final.
\newcommand{\lblPalInstruccionUno}{Przeczytaj słowo: najpierw sylabami, potem całe}
\newcommand{\lblPalInstruccionDos}{Przeczytaj oba słowa: najpierw sylabami, potem całe}
\newcommand{\lblPalInstruccionTres}{Przeczytaj trzy słowa: powoli, a potem płynnie}
\newcommand{\lblPalInstruccionCuatro}{Przeczytaj zdanie, pokazując palcem każde słowo}
\newcommand{\lblPalInstruccionSemana}{Słowa z tego tygodnia: przeczytaj je jeszcze raz}
\newcommand{\lblPalInstruccionSemanaFrases}{Zdania z tego tygodnia: przeczytaj je jeszcze raz}
\newcommand{\lblPalMarcaSemana}{Zaznacz każde, które już umiesz przeczytać.}
% Cuando lo único que se lee ese día es un nombre que se lee de un golpe
% (Lucía, Toby, Dani: ver PALABRAS_GLOBALES_PL en tools/gen_palabras.py).
\newcommand{\lblSlInstruccionNombre}{Przeczytaj to imię od razu, całe -- jak swoje własne}

\newcommand{\lblEscribe}{Napisz}
\newcommand{\lblEncuentra}{Znajdź}
\newcommand{\lblPalmadas}{Klaszcz}
\newcommand{\lblUne}{Połącz}
\newcommand{\lblSiNo}{Tak czy nie?}
\newcommand{\lblSi}{Tak}
\newcommand{\lblNo}{Nie}

% Instrucciones de actividad: las lee un adulto en voz alta -- por eso
% van en pequeño y en gris; lo que lee la niña o el niño va siempre
% grande.
\newcommand{\lblPalInstruccionCompleta}{Zamień te linie w rysunek.}
\newcommand{\lblPalInstruccionEscribe}{Obwiedź słowo ołówkiem, powoli, w środku liter. Potem napisz je samodzielnie na linii.}
\newcommand{\lblPalInstruccionEncuentra}{Zakreśl to słowo za każdym razem, kiedy je zobaczysz:}
\newcommand{\lblPalInstruccionPalmadas}{Przeczytaj każde słowo i klaśnij przy każdej sylabie. Pokoloruj jedno kółko za każde klaśnięcie.}
\newcommand{\lblPalInstruccionAdivina}{Dorosły czyta ci zagadkę. Narysuj odpowiedź.}
\newcommand{\lblPalInstruccionAdivinaEscribe}{Dorosły czyta ci zagadkę. Narysuj odpowiedź i napisz na linii, co to jest.}
\newcommand{\lblPalInstruccionUne}{Połącz linią każde słowo z tym samym słowem napisanym WIELKIMI LITERAMI.}
\newcommand{\lblPalInstruccionSiNo}{Przeczytaj każde zdanie. Zakreśl \lblSi\ albo \lblNo.}
\newcommand{\lblPalAhoraDibuja}{A teraz narysuj:}

\newcommand{\lblPalTituloPortada}{Uczę się czytać}
\newcommand{\lblPalNivelPortada}{Pierwsze słowa}
% En el cuaderno de verano (\segunEdicion, preamble.tex) ya no se lee una
% palabra al día, sino una frase corta.
\newcommand{\lblPalSubtituloPortada}{\segunEdicion{Codziennie nowe słowo -- i rysunek}{Codziennie krótkie zdanie -- i rysunek}}
\newcommand{\lblPalSubtituloPortadaMetodo}{\segunEdicion{Zeszyt do codziennego czytania dla tych, którzy zaczynają składać sylaby}{Zeszyt na lato dla tych, którzy czytają już słowa i zaczynają czytać zdania} -- krok przed zeszytem ze zdaniami}

% -----------------------------------------------------------------------
% El nivel 3 en polaco ("Leo con lupa", ver lang/es.tex): lo propio de
% ese nivel.
% -----------------------------------------------------------------------
\newcommand{\lblLupaInstruccionUno}{Czytaj na głos, bez pośpiechu. Zwróć uwagę na każdy szczegół}
\newcommand{\lblLupaInstruccionDos}{Przeczytaj dwa razy: najpierw dla siebie, a potem na głos}
\newcommand{\lblLupaInstruccionTres}{Czytaj na głos, głosem każdej postaci}
\newcommand{\lblLupaInstruccionCuatro}{Przeczytaj po cichu. Potem opowiedz to własnymi słowami}

\newcommand{\lblDibujaConLupa}{Narysuj z lupą}
\newcommand{\lblVinetas}{Komiks}
\newcommand{\lblMapa}{Mapa}
\newcommand{\lblLogica}{Tabela wskazówek}
\newcommand{\lblCaso}{Rozwiąż zagadkę}
\newcommand{\lblCodigo}{Tajna wiadomość}
\newcommand{\lblErrores}{Wyłap błędy}
\newcommand{\lblFicha}{Karta}
\newcommand{\lblCompara}{Porównaj}

\newcommand{\lblCompruebaDibujo}{Wróć do tekstu i sprawdź swój rysunek:}
\newcommand{\lblInstruccionCaso}{Zakreśl odpowiedź. Potem wypisz wskazówki z tekstu, które ci to mówią.}
\newcommand{\lblPistasCaso}{Wskazówki, które mi to mówią:}
\newcommand{\lblCorrigeErrores}{Napisz, jak jest naprawdę:}
\newcommand{\lblLupaInstruccionOrdena}[1]{Te zdarzenia są pomieszane. Wpisz w kratki #1, w takiej kolejności, w jakiej się wydarzyły.}
\newcommand{\lblLupaInstruccionVerdaderoFalso}{Zaznacz P (prawda) albo F (fałsz). Jeśli to fałsz, napisz pod spodem, jak jest naprawdę.}
\newcommand{\lblSalida}{START}
\newcommand{\lblLosDos}{oboje}

\newcommand{\lblLupaTituloPortada}{Czytam z lupą}
\newcommand{\lblLupaNivelPortada}{Uczę się czytać \textperiodcentered\ poziom 3}
\newcommand{\lblLupaSubtituloPortada}{Co tydzień nowa zagadka}
\newcommand{\lblLupaSubtituloPortadaMetodo}{Zeszyt do codziennego czytania dla tych, którzy czytają już płynnie: długie zdania i codziennie coś do odkrycia w tekście}

% -----------------------------------------------------------------------
% Las ediciones (ver lang/es.tex): el cuaderno de verano y la muestra
% gratuita.
% -----------------------------------------------------------------------
\newcommand{\lblInsigniaVerano}{Zeszyt na lato}
\newcommand{\lblMapaVeranoTitulo}{Lato, tydzień po tygodniu}
\newcommand{\lblMapaVeranoIntro}{Za każdy dzień czytania pokoloruj jedno słońce.}
\newcommand{\lblMapaDias}{Dni}
\newcommand{\lblMapaTema}{Co się dzieje}

% La muestra gratuita (las cuatro primeras semanas).
\newcommand{\lblInsigniaMuestra}{Bezpłatna próbka: pierwsze cztery tygodnie}
\newcommand{\lblMuestraFinTitulo}{Czytamy dalej?}
\newcommand{\lblMuestraFinTexto}[1]{To były pierwsze cztery tygodnie zeszytu \emph{#1}. Cały zeszyt ma \totaldias\ stron, po jednej na każdy dzień od poniedziałku do piątku przez cały rok, i można go pobrać za darmo, w kolorze albo w czerni i bieli, tutaj:}
